\section{Các thuật toán tìm kiếm}
\label{sec:algorithms}

Trên đồ thị hữu hạn của Chương~\ref{sec:model}, bốn thuật toán khác nhau ở cách
quản lý tập chờ xét (\emph{frontier}) và chọn trạng thái tiếp theo. Các mã giả
dưới đây tập trung vào lõi tìm kiếm. Trạng thái đầu được kiểm tra đích ngay;
các trạng thái còn lại được kiểm tra khi lấy ra khỏi tập chờ xét. Trạng thái
đích không được tính vào \texttt{expanded\_nodes} vì không sinh trạng thái kế
tiếp. Chi phí chuẩn đã được định nghĩa ở Bảng~\ref{tab:cost}; riêng BFS và DFS
coi mỗi hành động là một cạnh đơn vị.

Các bảo đảm đầy đủ và tối ưu dưới đây giả định quá trình tìm kiếm không bị
\texttt{CANCELLED}, \texttt{TIMEOUT} hoặc \texttt{NODE\_LIMIT}. Mã giả lược bỏ
các điều kiện dừng tùy chọn để tập trung vào lõi thuật toán. Các cận theo
$|V|,|E|$ giả định việc sinh trạng thái kế tiếp và tra cứu bảng băm có chi phí
đơn vị; khi kích thước màn chơi thay đổi, mỗi bước chuyển còn có thể duyệt
danh sách công tắc và cầu.

\subsection{Breadth-First Search}

BFS duyệt các trạng thái theo độ sâu tăng dần bằng hàng đợi FIFO. Bản cài đặt
dùng \texttt{deque} của thư viện \texttt{collections}. Từ điển \texttt{parent}
vừa lưu quan hệ để truy vết lời giải, vừa đánh dấu trạng thái đã thăm. Mỗi
trạng thái được đánh dấu khi đưa vào hàng đợi nên không bị thêm nhiều lần.

\begin{algorithm}[H]
\caption{BFS trong hệ thống}
\label{alg:bfs}
\begin{algorithmic}[1]
\Function{BFS}{$engine,start$}
\State $\mathit{frontier}\gets\langle start\rangle$ (FIFO)
\State $\mathit{parent}[start]\gets\varnothing$
\While{$\mathit{frontier}$ không rỗng}
    \State $s\gets\operatorname{popleft}(\mathit{frontier})$
    \If{$s$ là Goal} \State \Return \Call{Reconstruct}{$\mathit{parent},s$} \EndIf
    \State tăng $expanded$
    \ForAll{$(a,t)\in$ \texttt{engine.get\_successors(s)}}
        \If{$t\notin\mathit{parent}$}
            \State $\mathit{parent}[t]\gets(s,a)$
            \State $\operatorname{append}(\mathit{frontier},t)$
        \EndIf
    \EndFor
\EndWhile
\State \Return \textsc{NoSolution}
\EndFunction
\end{algorithmic}
\end{algorithm}

Trong BFS, mọi hành động có trọng số 1, kể cả \textsc{SwitchCube}. Do đó đường tìm
được tối ưu theo \texttt{action\_count}. Đại lượng
\texttt{solution\_length} chỉ đếm các hành động di chuyển, nên BFS không trực
tiếp tối ưu đại lượng này. \texttt{solution\_cost} của BFS bằng số hành động,
khác chi phí chuẩn của UCS và A*.

Với đồ thị hữu hạn, BFS đầy đủ và tối ưu theo số cạnh đơn vị. Bản tìm kiếm trên
đồ thị có thời gian $O(|V|+|E|)$ và bộ nhớ $O(|V|)$; theo ký hiệu cây tìm kiếm thường
gặp, số nút có thể tăng theo $O(b^d)$, trong đó $b$ là hệ số phân nhánh và $d$
là độ sâu lời giải nông nhất.

\subsection{Depth-First Search}

Nhóm chọn DFS tìm kiếm trên đồ thị, không dùng tìm kiếm sâu dần (IDS). Tập chờ xét
là \texttt{list} hoạt động như ngăn xếp. \texttt{GameEngine} trả các trạng thái
kế tiếp theo thứ tự Up, Down, Left, Right, SwitchCube; bản cài đặt đưa chúng
vào ngăn xếp theo thứ tự ngược để trạng thái đầu tiên được lấy ra trước.

\begin{algorithm}[H]
\caption{DFS có chống lặp}
\label{alg:dfs}
\begin{algorithmic}[1]
\Function{DFS}{$engine,start$}
\State $\mathit{frontier}\gets\langle start\rangle$ (stack)
\State $\mathit{parent}[start]\gets\varnothing$
\While{$\mathit{frontier}$ không rỗng}
    \State $s\gets\operatorname{pop}(\mathit{frontier})$
    \If{$s$ là Goal} \State \Return \Call{Reconstruct}{$\mathit{parent},s$} \EndIf
    \State tăng $expanded$
    \ForAll{$(a,t)$ theo thứ tự ngược của \texttt{engine.get\_successors(s)}}
        \If{$t\notin\mathit{parent}$}
            \State $\mathit{parent}[t]\gets(s,a)$
            \State $\operatorname{append}(\mathit{frontier},t)$
        \EndIf
    \EndFor
\EndWhile
\State \Return \textsc{NoSolution}
\EndFunction
\end{algorithmic}
\end{algorithm}

Tập trạng thái đã thăm ngăn DFS lặp vô hạn khi lăn ngược hoặc đổi khối điều
khiển hai lần. Trên đồ thị hữu hạn, DFS sẽ tìm thấy Goal nếu tồn tại đường đi
tới đó. Lời giải đầu tiên phụ thuộc vào thứ tự trạng thái kế tiếp và không bảo
đảm tối ưu số bước hay chi phí. Thời gian là $O(|V|+|E|)$; bộ nhớ của bản cài
đặt vẫn có thể đạt $O(|V|)$ vì phải lưu trạng thái đã thăm và quan hệ truy vết,
dù ngăn xếp DFS thường nhỏ hơn hàng đợi BFS.

\subsection{Uniform-Cost Search}

UCS áp dụng nguyên lý của Dijkstra trên đồ thị trạng thái~\cite{dijkstra1959}.
Mỗi phần tử trong hàng đợi ưu tiên chứa $(g,counter,state,path)$; biến
\texttt{counter} phân xử khi các giá trị ưu tiên bằng nhau mà không cần so sánh
hai trạng thái. Từ điển \texttt{best\_g} lưu chi phí nhỏ nhất đã biết. Phần
tử có $g$ lớn hơn giá trị hiện tại trong \texttt{best\_g} đã lỗi thời, nên bị
bỏ qua và không tính là một nút mở rộng.

UCS tối thiểu hóa tổng chi phí chuẩn, khác mục tiêu số hành động của BFS. Mọi chi phí ở
Bảng~\ref{tab:cost} đều không âm; \texttt{best\_g} chỉ cập nhật khi tìm được
đường có chi phí nhỏ hơn, nên \textsc{SwitchCube} với chi phí 0 không tạo vòng
lặp cập nhật vô hạn.

\begin{algorithm}[H]
\caption{Uniform-Cost Search}
\label{alg:ucs}
\begin{algorithmic}[1]
\Function{UCS}{$engine,start,cost$}
\State $\mathit{frontier}\gets\{(0,0,start,\langle\rangle)\}$ (min-heap)
\State $\mathit{bestG}[start]\gets0$; $\mathit{counter}\gets0$
\While{$\mathit{frontier}$ không rỗng}
    \State $(g,k,s,\mathit{path})\gets\operatorname{pop}(\mathit{frontier})$
    \If{$g>\mathit{bestG}[s]$} \State \textbf{continue} \EndIf
    \If{$s$ là Goal} \State \Return $(\mathit{path},g)$ \EndIf
    \State tăng $expanded$
    \ForAll{$(a,t)\in$ \texttt{engine.get\_successors(s)}}
        \State $\mathit{newG}\gets g+cost(s,a,t)$
        \If{$\mathit{newG}<\mathit{bestG}.\operatorname{get}(t,\infty)$}
            \State $\mathit{bestG}[t]\gets\mathit{newG}$
            \State $\mathit{counter}\gets\mathit{counter}+1$
            \State $\mathit{path}'\gets\mathit{path}\mathbin{\|}\langle a\rangle$
            \State $\operatorname{push}(\mathit{frontier},(\mathit{newG},\mathit{counter},t,\mathit{path}'))$
        \EndIf
    \EndFor
\EndWhile
\State \Return \textsc{NoSolution}
\EndFunction
\end{algorithmic}
\end{algorithm}

Trên đồ thị hữu hạn với chi phí cạnh không âm, UCS đầy đủ và tối ưu theo tổng
chi phí. Bản cài đặt sao chép chuỗi \texttt{path} khi đưa một đường tốt hơn
vào hàng đợi ưu tiên. Nếu $L$ là độ dài đường lưu lớn nhất, mỗi lần sao chép
có thể tốn $O(L)$; vì vậy bộ nhớ dành cho các phần tử chờ xét có thể tới
$O((|V|+|E|)L)$, còn thời gian có thêm phần sao chép bên cạnh thao tác hàng
đợi ưu tiên. UCS có thể chọn thêm hành động đổi khối với chi phí 0 hoặc tránh
ô Fragile có chi phí 2 để giảm tổng chi phí.

\subsection{A* Search}

A* chọn trạng thái theo giá trị ưu tiên
\begin{equation}
f(n)=g(n)+h(n),
\end{equation}
trong đó $g$ là chi phí từ trạng thái đầu, dùng cùng hàm chi phí với UCS, còn
$h$ ước lượng chi phí còn lại tới Goal. Heuristic \emph{chấp nhận được}
(admissible) nếu không vượt chi phí tối ưu còn lại; heuristic \emph{nhất quán}
(consistent) nếu ước lượng tại mỗi trạng thái không vượt chi phí của một cạnh
cộng ước lượng ở trạng thái kế tiếp. Phạm vi áp dụng của cận dưới $h$ được chứng
minh bên dưới.

Mỗi phần tử trong hàng đợi ưu tiên có dạng $(f,g,counter,state)$.
\texttt{best\_g} lưu đường rẻ nhất đã biết; khi tìm được giá trị $g$ nhỏ hơn,
A* có thể đưa cùng trạng thái vào hàng đợi lần nữa. Trạng thái đó có thể được
mở rộng lại, nên số lần mở rộng có thể lớn hơn số trạng thái phân biệt $|V|$.
\texttt{came\_from} lưu quan hệ truy vết mà không sao chép cả đường đi vào mỗi
phần tử. Cấu trúc này dựa trên A* tìm kiếm trên đồ thị của Hart, Nilsson và
Raphael~\cite{hart1968}.

\subsubsection{Heuristic cho khối nguyên}

Gọi $C(B)$ là tâm khối: $(r,c)$ khi đứng, $(r,c+0{,}5)$ khi nằm ngang và
$(r+0{,}5,c)$ khi nằm dọc. Với đường đi không kích hoạt Split Switch, mỗi
lần lăn dịch chuyển tâm theo khoảng cách Manhattan tối đa 1,5 và có chi phí
ít nhất 1. Vì thế, tổng chi phí để tâm tới ô đích $g$ không nhỏ hơn giá trị sau;
điều kiện phải đứng trên $g$ chỉ có thể làm số bước tăng thêm:
\begin{equation}
h_{normal}(B,g)=\frac{\lVert C(B)-g\rVert_1}{1{,}5}.
\end{equation}
Đây là cận dưới cho đường đi không dịch chuyển tức thời. Một lần kích hoạt
Split Switch có thể đổi vị trí khối vượt quá giới hạn dịch chuyển tâm nêu trên.

\subsubsection{Heuristic cho hai khối đơn vị}

Hai khối đơn vị ở $P_a,P_b$ phải kề trực giao để có thể ghép lại trước khi
khối nguyên hoàn thành màn chơi. Nếu khoảng cách giữa chúng là $d$, cần ít nhất
$\max(0,d-1)$ lần di chuyển một khối để đưa chúng đến vị trí ghép; đổi khối
đang điều khiển có chi phí 0 nhưng không đổi vị trí. Trên đường đi không dịch
chuyển tức thời, tổng khoảng cách Manhattan của hai khối tới đích giảm tối đa
3 trong một lần lăn khối nguyên và tối đa 1 trong một lần di chuyển khối đơn
vị. Mỗi lần di chuyển có chi phí ít nhất 1. Cả hai đại lượng đều là cận dưới,
nên ta lấy giá trị lớn hơn:
\begin{equation}
h_{split}(P_a,P_b,g)=\max\left(\max(0,\lVert P_a-P_b\rVert_1-1),
\frac{\lVert P_a-g\rVert_1+\lVert P_b-g\rVert_1}{3}\right).
\end{equation}

Trong bản cài đặt, $h_{direct}(s,x)$ áp dụng công thức tương ứng ở trên cho
khối nguyên hoặc hai khối đơn vị, với ô đích tạm $x$. Gọi $T_q$ là cặp vị trí
dịch chuyển của Split Switch $q$. Heuristic mặc định là
\begin{equation}
h(s)=\min\left(h_{direct}(s,g),1+\min_q h_{split}(T_q,g)\right),
\end{equation}
và bỏ nhánh thứ hai nếu màn chơi không có Split Switch. Đây là cận dưới theo
hàm chi phí chuẩn ở Bảng~\ref{tab:cost}. Xét một đường đi bất kỳ $\pi$ từ $s$
đến Goal. Nếu không có lần dịch chuyển tức thời nào, các cận hình học ở trên cho
$h_{direct}(s,g)\leq\operatorname{cost}(\pi)$. Nếu có dịch chuyển tức thời, xét
\emph{lần dịch chuyển cuối} tại $q$. Hành động kích hoạt này có chi phí ít nhất
1; đoạn còn lại không có dịch chuyển tức thời và có chi phí ít nhất
$h_{split}(T_q,g)$. Nếu hai vị trí đích kề nhau, $h_{split}$ bằng cận theo tâm
của khối vừa ghép vì Goal có tọa độ nguyên; việc ghép ngay vẫn giữ cận dưới.

Các hành động trước lần dịch chuyển tức thời cuối có chi phí không âm. Do đó
$h(s)\leq\operatorname{cost}(\pi)$ trong cả hai trường hợp, kể cả khi đường đi
dùng nhiều Split Switch. Ô Void, cầu đóng và ô Fragile chỉ loại bớt đường hoặc
tăng chi phí; chúng không làm mất cận dưới. Lấy đường tối ưu cho
$h(s)\leq h^*(s)$ trên mọi màn chơi hợp lệ với hàm chi phí chuẩn. Lập luận này
không chứng minh tính nhất quán của heuristic.

\begin{algorithm}[H]
\caption{A* với cập nhật đường tốt hơn}
\label{alg:astar}
\begin{algorithmic}[1]
\Function{AStar}{$engine,start,cost,h$}
\State $\mathit{frontier}\gets\{(h(start),0,0,start)\}$ (min-heap)
\State $\mathit{bestG}[start]\gets0$; $\mathit{cameFrom}\gets\varnothing$
\State $\mathit{counter}\gets0$
\While{$\mathit{frontier}$ không rỗng}
    \State $(f,g,k,s)\gets\operatorname{pop}(\mathit{frontier})$
    \If{$g>\mathit{bestG}[s]$} \State \textbf{continue} \EndIf
    \If{$s$ là Goal} \State \Return \Call{Reconstruct}{$\mathit{cameFrom},s$} \EndIf
    \State tăng $expanded$
    \ForAll{$(a,t)\in$ \texttt{engine.get\_successors(s)}}
        \State $\mathit{newG}\gets g+cost(s,a,t)$
        \If{$\mathit{newG}<\mathit{bestG}.\operatorname{get}(t,\infty)$}
            \State $\mathit{bestG}[t]\gets\mathit{newG}$
            \State $\mathit{cameFrom}[t]\gets(s,a)$
            \State $\mathit{counter}\gets\mathit{counter}+1$
            \State $f'\gets\mathit{newG}+h(t)$
            \State $\operatorname{push}(\mathit{frontier},(f',\mathit{newG},\mathit{counter},t))$
        \EndIf
    \EndFor
\EndWhile
\State \Return \textsc{NoSolution}
\EndFunction
\end{algorithmic}
\end{algorithm}

Khi $h=0$, A* trở thành UCS; nhóm dùng trường hợp này làm mốc kiểm thử. Với
heuristic mặc định và hàm chi phí chuẩn, A* trả về đường đi có tổng chi phí
nhỏ nhất trên đồ thị hữu hạn, kể cả khi có nhiều lần dịch chuyển tức thời. Kết
luận này đòi hỏi thuật toán được chạy đến khi lấy một trạng thái đích ra khỏi
hàng đợi, không bị hủy, hết thời gian hoặc chạm giới hạn nút. Bảo đảm này không
tự áp dụng cho hàm chi phí hoặc heuristic tùy chỉnh không thỏa điều kiện chấp
nhận được. Việc mở lại trạng thái còn có thể làm tăng thời gian và bộ nhớ; hiệu
quả thực tế được đánh giá ở Mục~\ref{sec:results}.
